\begin{dang}{HỆ YOUNG THAY ĐỔI CẤU TRÚC}
%\Anlythuyet{
%\Binhluan{
%\begin{itemize}
%\item Thay đổi bước sóng: $\dfrac{i_1}{i_2}=\dfrac{\lambda_1}{\lambda_2}$.
%\item Thay đổi khoảng cách giữa hai khe: $\dfrac{i_1}{i_2}=\dfrac{a_2}{a_1}$.
%\item Thay đổi khoảng cách từ hai khe tới màn: $\dfrac{i_1}{i_2}=\dfrac{D_1}{D_2}$
%\item Nếu thí nghiệm được tiến hành trong môi trường trong suốt có chiết suất n thì bước sóng và khoảng vân có mối liên hệ
%\begin{align*}
%\lambda_n=\dfrac{\lambda}{n}\Rightarrow i_n=\dfrac{\lambda_n.D}{a}=\dfrac{i}{n}
%\end{align*}
%\end{itemize}
%}{Nếu a tăng thì bậc vân giao thoa tại một điểm trên màn tăng và ngược lại.\\
%Nếu D hoặc $\lambda$ tăng thì bậc vân giao thoa tại một điểm trên màn giảm và ngược lại}
%}
\end{dang}